% Source: https://artofproblemsolving.com/community/c3h1744022_geo_questions
% Stable identity: aops-topic-1744022-q01
Use coordinates $0,1,2,3,4$ in each direction; the tunnels remove the unit cubes with at least two coordinates equal to $2$, i.e.\ $3\cdot5-2=13$ cubes, leaving $112$. After dipping, a face of a unit cube is painted when it is on the outside of the big cube or on the wall of a tunnel.

\begin{asy}
usepackage("amsmath");
size(4.5cm);
defaultpen(fontsize(10pt));
pair k=(0.5,0.35);
void face(pair o,pair u,pair v){
 for(int i=0;i<=5;++i){draw((o+i*u)--(o+i*u+5*v),linewidth(0.5));draw((o+i*v)--(o+i*v+5*u),linewidth(0.5));}
 fill((o+2*u+2*v)--(o+3*u+2*v)--(o+3*u+3*v)--(o+2*u+3*v)--cycle,black);
}
face((0,0),(1,0),(0,1));
face((0,5),(1,0),k);
face((5,0),k,(0,1));
draw((0,0)--(5,0)--(5,5)--(0,5)--cycle,linewidth(1));
draw((0,5)--(0,5)+5*k--(5,5)+5*k--(5,5),linewidth(1));draw((5,0)--(5,0)+5*k--(5,5)+5*k,linewidth(1));
label("\mbox{}",(min(currentpicture,user=true)+max(currentpicture,user=true))/2,black); // keeps the text after the figure black when LaTeX typesets the labels inline
\end{asy}

\begin{itemize}
\item \textbf{Edge cubes} (not corners) of the big cube: $12$ edges with $3$ each, $36$ cubes. Each has two outer faces and touches no tunnel, so exactly two painted faces.
\item \textbf{Face cubes} (one outer face): on each face, the four cubes sharing an edge with the hole have their outer face and one tunnel wall painted: $4\cdot6=24$ cubes. The other face cubes have one painted face.
\item \textbf{Interior cubes} (the $3\times3\times3$ core, minus the $7$ tunnel cubes in it): the $12$ cubes at the middles of the core's edges, such as $(2,1,1)$, touch two tunnels (at $(2,2,1)$ and $(2,1,2)$), so they have exactly two painted faces; the $8$ core corners such as $(1,1,1)$ touch no tunnel.
\item Corner cubes have three painted faces.
\end{itemize}
The total is $36+24+12=\boxed{72}$.
